\section*{Conventions de nommage}
On réutilise la structuration et les conventions de nommage suivante : 
\begin{description}
    \item[sections SFC] \texttt{<Mode>\_<NomDuMode>} ou \texttt{<Mode>\_<NomDuMode>\_Actions}
    \begin{itemize}
        \item Exemple : \texttt{F1\_ProductionNormale} pour le mode F1 complet ou \texttt{F1\_ProductionNormale\_Guichet} pour la partie guichet seule.
    \end{itemize}
    \item[transitions] \texttt{T\_<ModeSource>\_to\_<ModeCible>}
    \begin{itemize}
        \item Exemple : \texttt{T\_F1\_to\_F2} pour la transition du mode F1 vers le mode F2.
    \end{itemize}
    \item[macro-étapes] \texttt{<Mode>\_MS<n>\_<Action>}
    \begin{itemize}
        \item Exemple : \texttt{F1\_MS1\_PriseGobelet} pour la macro-étape 1 du mode F1, qui consiste à prendre un gobelet dans le guichet.
    \end{itemize}   
\end{description}


\section{Introduction}
\label{secIntroduction}

Ce document rassemble le \textbf{travail demandé} pour finaliser la mise en
service du poste guichet de la cellule robotique, en s'appuyant sur les
documents déjà établis :

\begin{itemize}
    \item \texttt{01\_conception} pour la présentation du système et le
    synoptique réseau ;
    \item \texttt{02\_analyse\_fonctionnelle} (mode~F1,
    \ref{secModeF1}) pour le comportement attendu de l'installation en
    production normale ;
    \item \texttt{02b\_donnees\_echangees} (\ref{secDonneesGuichet}) pour
    la description des capteurs et actionneurs du poste guichet gérés par
    le contrôleur PFC200 ;
    \item \texttt{02c\_configuration\_noc} pour la procédure générique de
    configuration d'un coupleur \texttt{BMX~NOC~0401.2} sous Control
    Expert.
\end{itemize}

\begin{UPSTIinfor}{Objet du document}
Durant ce TP, vous serez amenés à réaliser trois tâches principales :
\begin{enumerate}
    \item l'établissement de la \textbf{cartographie mémoire} \texttt{\%MW}
    des deux coupleurs NOC de l'installation, et la détermination des
    adresses exactes des capteurs et actionneurs du guichet ;
    \item l'application concrète de la procédure de configuration du
    coupleur \texttt{NOC\_ETHIP\_CONVOYEUR} sous Control Expert ;
    \item la programmation du mode~F1, \textbf{restreint au seul
    convoyeur/guichet} (sans le robot).
\end{enumerate}
\end{UPSTIinfor}
